\documentclass[11pt]{article}

\usepackage[margin=1in]{geometry}
\usepackage{amsmath, amssymb}
\usepackage{algorithm}
\usepackage{algpseudocode}
\usepackage{hyperref}
\usepackage{graphicx}

\title{Notes on Thompson Sampling}
\author{}
\date{}

\begin{document}
\maketitle

\section{The Multi-Armed Bandit Problem}

We have $d$ arms (in the ads example, $d = 10$ ads). At each round $n = 1, \dots, N$
we select one arm $i \in \{1, \dots, d\}$ and observe a reward
$r_i(n) \in \{0, 1\}$ (click / no click), drawn from an unknown Bernoulli
distribution with parameter $\theta_i$ specific to that arm:
\[
    r_i(n) \sim \text{Bernoulli}(\theta_i).
\]
The $\theta_i$ (true click-through rates) are unknown and must be estimated
online. The goal is to maximize cumulative reward, i.e. balance
\textbf{exploration} (trying arms we are unsure about) against
\textbf{exploitation} (playing the arm that currently looks best).

\section{Bayesian Idea}

Thompson Sampling treats each $\theta_i$ as a random variable and maintains a
posterior belief over it, updated after every observation. At each round it
\emph{samples} a value from each arm's posterior and picks the arm with the
highest sampled value --- i.e. it acts greedily with respect to a random draw
from what it currently believes, rather than a fixed statistic (like UCB's
upper confidence bound).

\subsection*{Prior}

Since rewards are Bernoulli, the conjugate prior for $\theta_i$ is the
Beta distribution. We start with an uninformative prior for every arm:
\[
    \theta_i \sim \text{Beta}(\alpha_i, \beta_i), \qquad
    \alpha_i = 1,\ \beta_i = 1 \quad (\text{i.e. Uniform}(0,1)).
\]

\subsection*{Posterior update (conjugacy)}

Let
\[
    N_i^1(n) = \text{number of times arm } i \text{ received reward } 1 \text{ up to round } n,
\]
\[
    N_i^0(n) = \text{number of times arm } i \text{ received reward } 0 \text{ up to round } n.
\]
Because Beta is conjugate to Bernoulli, the posterior after observing these
counts is again Beta:
\[
    \theta_i \mid N_i^1(n), N_i^0(n) \;\sim\; \text{Beta}\big(N_i^1(n) + 1,\; N_i^0(n) + 1\big).
\]
This is exactly \texttt{numbers\_of\_rewards\_1} and
\texttt{numbers\_of\_rewards\_0} in the implementation, and it is why the
sampling step is \verb|random.betavariate(N1[i] + 1, N0[i] + 1)|.

\section{Algorithm}

\begin{algorithm}[h]
\caption{Thompson Sampling for Bernoulli bandits}
\begin{algorithmic}[1]
\State $N_i^1 \gets 0,\; N_i^0 \gets 0$ for all arms $i = 1, \dots, d$
\For{$n = 1$ to $N$}
    \For{$i = 1$ to $d$}
        \State draw $\hat\theta_i \sim \text{Beta}(N_i^1 + 1,\, N_i^0 + 1)$
    \EndFor
    \State select arm $i^\ast = \arg\max_i \hat\theta_i$
    \State observe reward $r \in \{0, 1\}$ for arm $i^\ast$
    \If{$r = 1$}
        \State $N_{i^\ast}^1 \gets N_{i^\ast}^1 + 1$
    \Else
        \State $N_{i^\ast}^0 \gets N_{i^\ast}^0 + 1$
    \EndIf
\EndFor
\end{algorithmic}
\end{algorithm}

\section{Why It Works: Exploration vs. Exploitation}

\begin{itemize}
    \item Early on, with few observations, $\text{Beta}(1,1)$ or
        $\text{Beta}(\text{small}, \text{small})$ is wide/flat, so sampled
        $\hat\theta_i$ values vary a lot round to round $\Rightarrow$ arms get
        explored roughly uniformly.
    \item As $N_i^1, N_i^0$ grow for a given arm, the Beta posterior
        concentrates around the empirical rate
        $\frac{N_i^1}{N_i^1 + N_i^0}$, so sampled values become
        consistent and predictable $\Rightarrow$ the algorithm exploits the
        arm if it truly is good.
    \item Arms that look bad get sampled less, so their posteriors stay
        wide for longer, keeping a nonzero (but shrinking) chance they still
        get explored, in case early evidence was misleading.
\end{itemize}

This randomized-but-principled trade-off is why Thompson Sampling tends to
outperform simpler heuristics and even Upper Confidence Bound (UCB) in
practice, especially with delayed/batched feedback.

\section{Thompson Sampling vs. UCB}

\begin{center}
\begin{tabular}{|l|l|l|}
\hline
& \textbf{UCB} & \textbf{Thompson Sampling} \\
\hline
Nature & Deterministic & Probabilistic (stochastic) \\
\hline
Selection rule & $\arg\max_i \left(\bar r_i + \sqrt{\tfrac{3}{2}\tfrac{\log n}{n_i}}\right)$
    & $\arg\max_i \hat\theta_i,\ \hat\theta_i \sim \text{Beta}(N_i^1+1, N_i^0+1)$ \\
\hline
Feedback timing & Needs update every round & Can tolerate delayed/batched updates \\
\hline
Empirical performance & Good & Often better in practice \\
\hline
\end{tabular}
\end{center}

\section{When to Use Which}

\begin{itemize}
    \item \textbf{UCB} is deterministic and easy to reason about: given the
        same history, it always makes the same choice. It needs the full
        cumulative history recomputed every round, and it forces every arm
        to be tried once up front (an infinite/very large upper bound is
        assigned to unseen arms, seen in code as
        \texttt{upper\_bound = 1e400}). It has a tunable exploration
        constant (the $\sqrt{3/2 \cdot \cdot}$ term) that controls how
        aggressively it explores. Good fit for stationary problems where
        reward feedback arrives immediately every round and you want a
        reproducible, auditable selection rule.
    \item \textbf{Thompson Sampling} is probabilistic: it only needs the
        latest success/failure counts (not the full history) to draw a
        sample, so it updates cheaply and incrementally. There is no
        tunable exploration parameter --- the Beta(1,1) prior is naturally
        wide for unseen arms, so cold start is handled for free. It also
        tends to tolerate delayed or batched reward feedback better, since
        stale posteriors just stay wide instead of stalling the algorithm.
        This is why Thompson Sampling is generally preferred in production
        settings such as ad serving, online A/B testing, and
        recommendation systems, and why it converges faster / earns a
        higher \texttt{total\_reward} than UCB on this dataset in practice.
\end{itemize}

\section{Correspondence to the Notebook Code}

\begin{center}
\begin{tabular}{|l|l|}
\hline
\textbf{Math} & \textbf{Code (\texttt{thompson\_sampling.ipynb})} \\
\hline
$d$ & \texttt{d = 10} (number of ads) \\
$N$ & \texttt{N = 10000} (number of rounds/users) \\
$N_i^1(n)$ & \texttt{numbers\_of\_rewards\_1[i]} \\
$N_i^0(n)$ & \texttt{numbers\_of\_rewards\_0[i]} \\
$\hat\theta_i \sim \text{Beta}(N_i^1+1, N_i^0+1)$ & \texttt{random.betavariate(N1[i]+1, N0[i]+1)} \\
$i^\ast = \arg\max_i \hat\theta_i$ & \texttt{if random\_beta > max\_random: ad = i} \\
cumulative reward & \texttt{total\_reward} \\
\hline
\end{tabular}
\end{center}

\end{document}
